\section{屈折・反射・表面波：地下を拘束する情報の違い}
\label{theory:exploration:seismic}
地震探査の記録には直接波、屈折波、反射波、表面波が重なる。屈折法では通常初動走時、反射法では反射イベントの走時と波形、表面波法では周波数ごとの位相速度を主に使う。名前の違いは処理だけではなく、地下構造を拘束する観測情報の違いである。

\subsection{速度と地質：波種・状態・境界を分ける}
\label{theory:exploration:waveproperties}
安定な等方弾性体では、体積弾性率 \(K\)、せん断弾性率 \(G\)、密度 \(\rho\) を使い、
\[
 V_P=\sqrt{\frac{K+4G/3}{\rho}},\qquad
 V_S=\sqrt{\frac{G}{\rho}}
\]
となる。\(K>0,G>0\) なら \(V_P>V_S\) である。流体ではせん断を保てず通常のS波が伝わらないので、二つの伝播速度を常に観測して比較できるとは限らない。

岩種が同じでも、割れ目の開口・接触、風化、変質、空隙、飽和が有効弾性係数を変える。空隙の空気を水が置き換えると圧縮性等へ影響し、P波速度は変わり得る。USGSの測定報告では飽和帯の深さとP波速度を読み、表面波のS波逆解析を拘束する\cite{usgs_refraction2013}。地質学の境界は岩種・年代・堆積単位等の区分であり、速度差の境界と同一ではない。速度から風化や破砕を推定するとき、露頭やボーリングと併せて解釈する。

通常の浅部屈折ではP波が多く使われるが、S波屈折も存在する\cite{epa_seismic_refraction}。土質のS波調査を比較的平坦な地形で行うという運用の目安を、急峻地形では物理的に不可能という命題へ変更しない。強震動にはS波構造が重要であり、通常の地山P波屈折の主用途と、複数手法の統合による寄与を分ける\cite{segj_integrated2017,japansea_spac2017}。

\subsection{屈折初動：はぎとりと往復の整合}
\label{theory:exploration:refraction}
上層速度 \(V_1\)、下層速度 \(V_2>V_1\)、上層厚 \(h\) の単純な水平二層では、\(\sin i_c=V_1/V_2\) を満たす臨界屈折波の走時は
\[
 t_{\rm head}(x)=t_i+\frac{x}{V_2},\qquad
 t_i=\frac{2h\cos i_c}{V_1}
\]
となる。直接波の \(x/V_1\) と比べて最初に到達する時刻を読む。原点走時 \(t_i\) は走時直線の外挿切片であり、起振位置の実測初動ではない。臨界距離と直接波を追い越す交差距離も同じ概念ではない\cite{epa_seismic_refraction}。

萩原のはぎとり法は表層の走時寄与を分離して下層の速度と境界を求める屈折初動の方法である。JICA調査団の比較表は直接波＋屈折波を扱い、トモグラフィとの構造・速度の仮定を分けている\cite{jica_refraction2015}。反射波の走時をそのまま同じはぎとりへ入力する説明は誤りである。

端点AからBへの全走時とBからAへの全走時は、相反的な同じ経路・媒体で整合する。全走時一致はこの往復整合などを確かめる意味であり、全ての受振点の到達時刻が同じという意味ではない。原点走時の整合、曲線の平行性も観測モデルに沿って確認する。薄い中間層は屈折波があっても初動として現れない距離範囲しか持たない場合がある。低速度逆転、薄いブラインド層、連続速度勾配を分ける\cite{usgs_refraction2013}。

連続的に深度増速するミラージュ構造では波線が曲がり、走時は折れ線ではなく曲線となる。トモグラフィは波線が通る格子速度又は遅さを反復更新し、連続勾配も扱える。ただし通らない領域の速度を一意に決めることはできない。測線は地形と地質の変化を見ながら設計し、等高線に垂直で一定勾配という望ましい典型配置だけを絶対条件にしない。

\subsection{反射走時：双曲線と放物線の適用範囲}
\label{theory:exploration:reflectiongeometry}
単一水平反射面、上部の一定速度 \(V\)、深さ \(h\)、震源と受振器の距離 \(x\) なら、反射点は両者の中点にある。左右対称の経路長を足して独自に導くと、
\[
 t(x)=\frac{2}{V}\sqrt{h^2+(x/2)^2},\qquad
 t^2=t_0^2+\frac{x^2}{V^2},\quad t_0=\frac{2h}{V}.
\]
\(t\)-\(x\) 平面では双曲線である。多層では小〜中オフセットの範囲で有効な重合速度を使って双曲線近似する。大オフセット、強い不均質、異方性等には非双曲線成分が残る。

さらに \(|x|\ll Vt_0\) なら平方根を展開して、
\[
 t(x)=t_0+\frac{x^2}{2V^2t_0}
 -\frac{x^4}{8V^4t_0^3}+\cdots
\]
となる。二次項までの近似は放物線である。このため平成28年度問21と令和5年度問32の公式選択を、放物線ではどんな近似もできないという命題へ書き換えない。通常の走時モデルが双曲線であることと、原文に近似範囲の指定がない問題点を併記する。

NMOは各オフセットの走時をゼロオフセット走時へそろえる。重合速度はその目的に最適な速度であり、真の区間速度と無条件に同じではない\cite{usgs_processing2000}。反射とは異なるmoveoutの多重反射等は重合で弱まる場合があり、相関雑音も全て残るとは限らない。同じmoveoutの雑音は抑えにくい。

浅層反射の最大オフセットを目標深度程度とする説明は通常の展開目安である\cite{coastal_reflection2016}。目標反射角、近接雑音、速度解析、深度、開口と重合数で設計を変える。二倍を一律に必要とする説明を誤りと選ぶ場合も、二倍の観測が常に不適切という意味ではない。

\subsection{分解能と検出：薄層が存在することを知る}
\label{theory:exploration:resolution}
厚さ \(\Delta h\) の上下面反射の時間差は単純な垂直経路で \(\Delta t=2\Delta h/V\) である。二つのwaveletが強く重なると別々の反射として分離しにくくなる。代表的な垂直分解能の目安 \(\Delta h\sim\lambda/4\) は、上面と下面の分離に関するものである\cite{seg_vertical_resolution}。

分離できない薄層も、合成振幅や周波数応答の変化として存在を検出できる場合がある。検出にはS/N、インピーダンス差、位相、帯域が関わる。「四分の一波長より薄ければ必ず見えない」と断定しない。Rayleigh基準やRicker基準、waveletは反射イベントの分離に関係し、屈折直線のintercept timeとは役割が異なる。

横方向にはFresnel zoneの有限範囲から反射が寄与する\cite{seg_lateral_resolution}。単純なゼロオフセットの幾何では往復経路差を半波長とする一次zoneから \(r_F\approx\sqrt{h\lambda/2}\) の目安を独自に得る。migration、開口、帯域と速度モデルによって横方向の分解能は変わる。zone径を全処理の最終解像度と固定しない。

\subsection{表面波分散とS波モデル}
\label{theory:exploration:surfacewaves}
Rayleigh波は表面に沿って進み、深度へ振幅が減少する。不均質な成層では周波数により感度深度が変わるため位相速度 \(c(f)\) が分散する。長波長は深部、短波長は浅部を拘束する傾向があるが、一つの波長が一つの深度だけを測るわけではない\cite{epa_surfacewaves}。均質等方半無限体は非分散という退化例である。

観測位相速度から地下モデルを作り、理論分散曲線を計算して残差を小さくする。位相速度は \(V_P,V_S,\rho\)、層厚に依存するが、通常の地盤分散は \(V_S\) に強い感度を持つ。Xiaらの一次論文要約は高周波側でS波速度と層厚の影響が大きいことを述べる\cite{xia_rayleigh1999}。P波・密度の影響をゼロにせず、別の測定や拘束を使う。

水平多層の一次元逆解析は通常のモデルであり、地盤自体が常に横方向に均質と分かったわけではない。高次モード、方向による変化、モデル自由度とデータ数を確認する。能動表面波の通常の浅部対象を20m程度とする資料がある\cite{moe_surface2004}。しかし20mや100mは普遍的上限ではなく、震源の長波長、受振器の低周波応答、配置、雑音と解析範囲に依存する。微動で長波長側を補えれば深部拘束を増やせるが、併用だけで深度増加を保証しない。

\subsection{微動アレイの短波長と長波長の限界}
\label{theory:exploration:arraylimits}
空間標本化で間隔 \(d\) に対し波長が短すぎると、異なる空間波数が同じ標本を与える。等間隔・単方向の単純モデルでは \(\lambda>2d\) の条件が基本となる。実アレイの限界は配置と波の方向にも依存する。空間aliasingは主に最短波長側の制約であり、最長側と取り違えない\cite{japansea_spac2017}。

長波長ではアレイ内の位相差が小さく、有限開口、ノイズ、推定算法と観測時間の制約が効く。F-Kは周波数ごとの波数分布を求め、SPACは空間自己相関と分散の関係を使う。通常の条件でSPACがF-Kより長波長まで推定する例が一次論文にあるが、入射方向と多数の波、構造、配置、実装に依存する\cite{array_limits2002}。一定の半径倍率を全観測の保証値へ転記しない。複数サイズで重なる帯域を確認し、疑わしい分散曲線を機械的に一つのモデルへ合わせない。
